\documentclass[a4paper,9pt]{extarticle}

%=========================================================
% PACKAGES
%=========================================================

\usepackage[
a4paper,
top=1.8cm,
bottom=1.5cm,
footskip=7mm,
left=1cm,
right=1cm,
headheight=14pt,
headsep=5mm
]{geometry}
\usepackage[T1]{fontenc}
\usepackage[utf8]{inputenc}
\usepackage[french]{babel}

\usepackage[sfdefault]{FiraSans}
\renewcommand{\familydefault}{\sfdefault}

\usepackage{multicol}
\setlength{\columnsep}{0.7cm}

\usepackage{amsmath,amssymb}

\usepackage{titlesec}
\usepackage{enumitem}

\usepackage{tcolorbox}
\tcbuselibrary{skins,breakable}

\usepackage{xcolor}

\usepackage{graphicx}
\usepackage{float}

\usepackage{tikz}

\usepackage{fancyhdr}
\usepackage[hidelinks]{hyperref}


%=========================================================
% PAGE SETTINGS
%=========================================================

\setlength{\parindent}{0pt}
\setlength{\parskip}{1.5pt}

\setlength{\abovedisplayskip}{4pt}
\setlength{\belowdisplayskip}{4pt}

\setlist[itemize]{
leftmargin=*,
itemsep=1pt,
topsep=1pt
}


%=========================================================
% COLORS
%=========================================================

\definecolor{Blue}{HTML}{E3F0FF}
\definecolor{Green}{HTML}{E3F7EA}
\definecolor{Pink}{HTML}{FFE5ED}
\definecolor{Orange}{HTML}{FFF0DC}
\definecolor{Purple}{HTML}{EEE7FF}
\definecolor{Yellow}{HTML}{FFFADB}


\definecolor{BlueTitle}{HTML}{285B8F}
\definecolor{GreenTitle}{HTML}{287052}
\definecolor{PinkTitle}{HTML}{A84063}
\definecolor{OrangeTitle}{HTML}{A96500}
\definecolor{PurpleTitle}{HTML}{6245A5}


%=========================================================
% HEADER
%=========================================================

\pagestyle{fancy}

\fancyhf{}

\fancyhead[L]{\small\textbf{Game Theory and Strategic Decisions}}
\fancyhead[C]{\small Dominance and Iterated Elimination}
\fancyhead[R]{\small Week 02 --- 14-20/09/2026}

\fancyfoot[C]{\small\href{https://github.com/AlphaXT38G/notes_ba5}{github.com/AlphaXT38G/notes\_ba5}}
\fancyfoot[R]{\small \thepage}

\renewcommand{\headrulewidth}{0.4pt}



%=========================================================
% TITLES
%=========================================================

\titleformat{\section}
{\large\bfseries\color{BlueTitle}}
{}{0pt}{}

\titlespacing{\section}
{0pt}{6pt}{3pt}


% separator

\newcommand{\separator}{
\vspace{2pt}
\hrule
\vspace{3pt}
}


%=========================================================
% BOXES
%=========================================================

% Boxes are kept whole (unbreakable): inside multicols, a box split by tcolorbox
% leaves an empty column behind it.

% DEFINITION - blue
\newtcolorbox{definition}{
enhanced,
unbreakable,
colback=Blue,
colframe=BlueTitle,
boxrule=0.7pt,
arc=0mm,
title=\textbf{Definition},
fonttitle=\bfseries
}

% IMPORTANT - pink
\newtcolorbox{important}{
enhanced,
unbreakable,
colback=Pink,
colframe=PinkTitle,
boxrule=0.7pt,
arc=0mm,
title=\textbf{Key idea},
fonttitle=\bfseries
}

% FORMULAS - green
\newtcolorbox{formula}{
enhanced,
unbreakable,
colback=Green,
colframe=GreenTitle,
boxrule=0.7pt,
arc=0mm,
title=\textbf{Key formulas},
fonttitle=\bfseries
}

% WARNING - orange
\newtcolorbox{warning}{
enhanced,
unbreakable,
colback=Orange,
colframe=OrangeTitle,
boxrule=0.7pt,
arc=0mm,
title=\textbf{Warning},
fonttitle=\bfseries
}

% EXAMPLE - purple
\newtcolorbox{example}{
enhanced,
unbreakable,
colback=Purple,
colframe=PurpleTitle,
boxrule=0.7pt,
arc=0mm,
title=\textbf{Example},
fonttitle=\bfseries
}

% SUMMARY - orange
\newtcolorbox{summary}{
enhanced,
unbreakable,
colback=Yellow,
colframe=OrangeTitle,
boxrule=0.7pt,
arc=0mm,
title=\textbf{Summary},
fonttitle=\bfseries
}



%=========================================================
% FIGURE BOX
%=========================================================

\newtcolorbox{figurebox}{
enhanced,
colback=white,
colframe=gray,
boxrule=0.5pt,
arc=0mm,
center
}



%=========================================================
% DOCUMENT
%=========================================================

\begin{document}

%\footnotesize   --> To reduce the size of the text

\thispagestyle{fancy}



\begin{multicols}{2}



%=========================================================
\section{Recap: the Battle of the Sexes}


This game summarizes the techniques of the previous lecture: how to find a pure-strategy NE,
how to find a mixed-strategy NE, and that the same game can have both types of NE.


\begin{example}

\textbf{Story}

Two friends prefer to go out together than to go out alone. But Adam prefers to go to a
soccer game and Bob prefers to go to the opera. They don't have phones so cannot
communicate before deciding where to go.

\begin{center}
\begin{tabular}{c|cc}
Adam $\backslash$ Bob & Soccer & Opera \\
\hline
Soccer & (2,1) & (0,0) \\
Opera & (0,0) & (1,2) \\
\end{tabular}
\end{center}

\end{example}


\textbf{Pure-strategy NE}

\begin{itemize}

\item If Adam goes to soccer, also Bob prefers to go to soccer. If Adam goes to the opera,
also Bob prefers to go to the opera.

\item If Bob goes to soccer, also Adam prefers to go to soccer. If Bob goes to the opera,
also Adam prefers to go to the opera.

\end{itemize}


\begin{important}

\textbf{Soccer-Soccer} and \textbf{Opera-Opera} are the two pure-strategy NE.

\end{important}



%=========================================================
\section{Mixed-Strategy Nash Equilibria}


\begin{important}

In a game with $N$ actions, to adopt a mixed strategy
$s^{(i)} = (p^{(1)}, p^{(2)}, \dots, p^{(N)})$, a player must be \textbf{indifferent
between all the actions to which the strategy assigns non-zero probability}.

\textbf{Why?} Because if one of these actions gave lower expected payout than the other
actions, the player would never adopt it.

\end{important}


This is the principle that guides us to find mixed-strategy Nash Equilibria.


\begin{formula}

\textbf{Bob indifferent}: Adam plays $s^{(A)} = (p_S^{(A)}, 1 - p_S^{(A)})$.

\begin{itemize}

\item Bob goes to soccer:
$p_S^{(A)} \times 1 + (1 - p_S^{(A)}) \times 0 = p_S^{(A)}$

\item Bob goes to the opera:
$p_S^{(A)} \times 0 + (1 - p_S^{(A)}) \times 2 = 2 - 2p_S^{(A)}$

\end{itemize}

\[
\Rightarrow \quad p_S^{(A)} = 2/3, \quad s^{(A)} = (2/3, 1/3)
\]

\textbf{Adam indifferent}: Bob plays $s^{(B)} = (p_S^{(B)}, 1 - p_S^{(B)})$.

\begin{itemize}

\item Adam goes to soccer:
$p_S^{(B)} \times 2 + (1 - p_S^{(B)}) \times 0 = 2p_S^{(B)}$

\item Adam goes to the opera:
$p_S^{(B)} \times 0 + (1 - p_S^{(B)}) \times 1 = 1 - p_S^{(B)}$

\end{itemize}

\[
\Rightarrow \quad p_S^{(B)} = 1/3, \quad s^{(B)} = (1/3, 2/3)
\]

\end{formula}


\begin{important}

\textbf{Mixed-strategy Nash equilibrium}:
\[
s^{(A)} = (2/3, 1/3), \qquad s^{(B)} = (1/3, 2/3)
\]

\end{important}



\end{multicols}


\end{document}
